% ==================================================================
% id:       sets-003
% topic:    תורת הקבוצות ומאורעות
% subtopic: פעולות על קבוצות ומאורעות
% solution: full
% ==================================================================
יהיו $A,B,C$ קבוצות. הוכיחו בעזרת הכלה דו-כיוונית:
\begin{enumerate}
\item $(A\cup B)^C=A^C\cap B^C$.
\begin{solbox}
נוכיח שתי הכלות.

\textbf{$(A\cup B)^C\subseteq A^C\cap B^C$:} יהי $x\in(A\cup B)^C$. אז $x\notin A\cup B$, כלומר $x$ אינו באף אחת משתי הקבוצות: $x\notin A$ וגם $x\notin B$. לכן $x\in A^C$ וגם $x\in B^C$, כלומר $x\in A^C\cap B^C$.

\textbf{$A^C\cap B^C\subseteq(A\cup B)^C$:} יהי $x\in A^C\cap B^C$. אז $x\in A^C$ וגם $x\in B^C$, כלומר $x\notin A$ וגם $x\notin B$. לכן $x$ אינו באף אחת משתי הקבוצות, כלומר $x\notin A\cup B$, ולכן $x\in(A\cup B)^C$.

משתי ההכלות נובע ש-$\ans{(A\cup B)^C=A^C\cap B^C}$.
\end{solbox}
\item $A\cap(B\cup C)=(A\cap B)\cup(A\cap C)$. ציירו דיאגרמת ון לכל אחד מהאגפים.
\begin{solbox}
נוכיח שתי הכלות.

\textbf{$A\cap(B\cup C)\subseteq(A\cap B)\cup(A\cap C)$:} יהי $x\in A\cap(B\cup C)$. אז $x\in A$ וגם $x\in B\cup C$, כלומר $x\in B$ או $x\in C$. אם $x\in B$, אז $x\in A\cap B$; אם $x\in C$, אז $x\in A\cap C$. בשני המקרים $x\in(A\cap B)\cup(A\cap C)$.

\textbf{$(A\cap B)\cup(A\cap C)\subseteq A\cap(B\cup C)$:} יהי $x\in(A\cap B)\cup(A\cap C)$. אם $x\in A\cap B$, אז $x\in A$ ו-$x\in B\subseteq B\cup C$; אם $x\in A\cap C$, אז $x\in A$ ו-$x\in C\subseteq B\cup C$. בשני המקרים $x\in A$ וגם $x\in B\cup C$, כלומר $x\in A\cap(B\cup C)$.

לכן $\ans{A\cap(B\cup C)=(A\cap B)\cup(A\cap C)}$. בשני האגפים האזור הצבוע הוא אותו אזור: החלק של $A$ שנמצא ב-$B$ או ב-$C$.

\begin{center}
\venn{1/1/0,1/0/1,1/1/1}{$A\cap(B\cup C)$}\qquad
\venn{1/1/0,1/0/1,1/1/1}{$(A\cap B)\cup(A\cap C)$}
\end{center}
\end{solbox}
\end{enumerate}
